\subsection{灵敏度分析：重新估计后的厚度变化}
\label{sec:sensitivity-analysis}

问题一先固定同一组合成光谱，考察光学条件改变后厚度怎样移动，再将这一比较延伸至实测材料。每项扰动均重新估计厚度，其余未固定参数随之重估。

合成例在折射率、色散斜率、偏振、平滑和角度扰动下，厚度落在$6.66$--$10.01\,\mu\mathrm{m}$，最大绝对漂移为$2.01\,\mu\mathrm{m}$。双束界面场反演仍给出有限厚度，但未知光学条件足以使它发生微米量级移动。
% src:求解/问题1/结果/灵敏度.json:条件压力范围_微米;最大绝对漂移_微米;重估结果

实测碳化硅放大了厚度与折射率的互相补偿：切分、剖面和满足参数约束的情景合并后，厚度范围为$0.64$--$18.11\,\mu\mathrm{m}$；五折极差除以中位数为$0.7604$，无量纲。图\ref{fig:q2-thickness}分开呈现这些来源，问题三的碳化硅沿用同一范围。
% src:求解/问题2/结果/厚度结果.json:条件区间_微米;条件区间构造
% src:求解/问题2/结果/验证.json:五折跨切分稳定性.相对极差
% src:求解/问题3/结果/回炉轮3候选/仲裁闭环/20260910_131501_81623_建模公式/厚度结果.json:核心指标.碳化硅正式条件范围下限_um;核心指标.碳化硅正式条件范围上限_um

为分开光学变化与优化继续推进的影响，硅次折增加同起点的未扰动控制：控制厚度为$16.0477\,\mu\mathrm{m}$，衬底对比增加$20\%$后为$13.8276\,\mu\mathrm{m}$，扣除控制后的厚度变化为$-2.2201\,\mu\mathrm{m}$。两个终点均满足收敛标准，因而这对差值对应同一起点附近的光学条件响应。
% src:求解/问题3/结果/解读核验_回炉3_仲裁后.json:灵敏度.示例控制厚度_um;灵敏度.示例扰动厚度_um;灵敏度.示例扣控制变化_um;灵敏度.示例成对收敛
% src:求解/问题3/结果/回炉轮3候选/仲裁闭环/20260910_131501_81623_建模公式/厚度结果.json:同起点控制分解.0.情景
% src:交接/实验记录.json:1724.尝试;1724.现象;1724.决定

原第三版的参数扰动保留为历史对照：硅厚度变化$-0.4571$至$3.9259\,\mu\mathrm{m}$，碳化硅变化$-1.6720$至$2.1261\,\mu\mathrm{m}$。图\ref{fig:19}和表\ref{tab:10}保留这些相对原折基准的差值；它们没有扣除未扰动优化推进，与同起点控制差分别解释。
% src:求解/问题3/结果/灵敏度.json:灵敏度_参数扰动.29.厚度变化_um;灵敏度_参数扰动.27.厚度变化_um;灵敏度_参数扰动.73.厚度变化_um;灵敏度_参数扰动.72.厚度变化_um

\begingroup
\small
\setlength{\LTpre}{0.3em}
\setlength{\LTpost}{0.3em}
\begin{longtable}{>{\centering\arraybackslash}p{0.14\textwidth}>{\centering\arraybackslash}p{0.235\textwidth}>{\centering\arraybackslash}p{0.19\textwidth}>{\centering\arraybackslash}p{0.17\textwidth}>{\centering\arraybackslash}p{0.105\textwidth}}
\caption{扰动重估改变厚度与拟合损失}\label{tab:10}\\
\toprule
重估对象 & 扰动设置 & 厚度或差/$\mu\mathrm{m}$ & 训练损失 & 判断\\
\midrule
\endfirsthead
\toprule
重估对象 & 扰动设置 & 厚度或差/$\mu\mathrm{m}$ & 训练损失 & 判断\\
\midrule
\endhead
合成厚度 & 折射率曲线$\pm20\%$ & $6.66$--$10.01$ & --- & 分支移动\\
碳化硅厚度 & 原设置 & $10.80$ & $0.2747$ & 参照\\
碳化硅厚度 & 折射率$-3\%/+3\%$ & $8.21/7.96$ & $0.3384/0.3294$ & 分支移动\\
碳化硅厚度 & 窗口下/上界$+100$ & $11.41/8.89$ & $0.2005/0.3161$ & 分支移动\\
碳化硅厚度 & 网格$96\to192$点 & $13.53$ & $0.2744$ & 分支移动\\
硅历史差 & 光学、响应与窗口 & $-0.4571$--$3.9259$ & --- & 分支移动\\
碳化硅历史差 & 光学、响应与窗口 & $-1.6720$--$2.1261$ & --- & 分支移动\\
\bottomrule
\end{longtable}
% src:求解/问题1/结果/灵敏度.json:重估结果.1.缩放;重估结果.2.缩放;条件压力范围_微米
% src:求解/问题2/结果/验证.json:窗口与模型灵敏度.0;窗口与模型灵敏度.5;窗口与模型灵敏度.6;窗口与模型灵敏度.8;窗口与模型灵敏度.9;窗口与模型灵敏度.10
% src:求解/问题3/结果/灵敏度.json:灵敏度_参数扰动.29.厚度变化_um;灵敏度_参数扰动.27.厚度变化_um;灵敏度_参数扰动.73.厚度变化_um;灵敏度_参数扰动.72.厚度变化_um
% src:交接/实验记录.json:564.现象
\endgroup

网格加密使碳化硅厚度移动$25.27\%$，而原代表点的损失仅比补查最低值高$0.1167\%$。这是固定点集、模型和物理边界后的数值搜索变化，不能与改变光学假设的情景变化混称为同一种不确定性。表\ref{tab:q2-search-check}保留两类近优点及其边界状态。
% src:求解/问题2/结果/搜索补查.json:原网格到加密厚度变化_百分比;原代表相对最低损失差_百分比

表\ref{tab:10}中窗口移动以$\mathrm{cm}^{-1}$计，斜线前后逐项对应；不同窗口的损失对应不同点集。异常原值均在主拟合窗口外，保留原值并不改变拟合点集，故裁剪异常不能解释厚度改善。
% src:求解/问题2/结果/验证.json:窗口与模型灵敏度.0;窗口与模型灵敏度.8
% src:求解/问题2/结果/数据审查.json:质量处置

窗口比较曾沿用原波段，两种边界结果重合；因此改为各向内移动$100\,\mathrm{cm}^{-1}$，在既定训练成员与新窗口的交集上重估。两折按所移动的边界分别保留每角$273$或$292$个训练点，选段影响才得以分开。
% src:交接/实验记录.json:554.尝试;554.现象;555.现象;556.尝试
% src:求解/问题3/结果/灵敏度.json:灵敏度_参数扰动.15.请求固定参数值;灵敏度_参数扰动.16.请求固定参数值;灵敏度_参数扰动.15.训练源行;灵敏度_参数扰动.16.训练源行;灵敏度_参数扰动.36.训练源行;灵敏度_参数扰动.37.训练源行;灵敏度_参数扰动.57.训练源行;灵敏度_参数扰动.58.训练源行;灵敏度_参数扰动.78.训练源行;灵敏度_参数扰动.79.训练源行

光学与角度扰动取折射率$\pm3\%$、衬底对比及损耗$\pm20\%$、入射角$\pm0.5^\circ$，偏振分别取两种纯态。重估后，硅沿用完整往返场，碳化硅保留两束基准。厚度响应记录这些设置对估计的影响，窗口和网格比较则揭示分支移动；下一章据此评价方法与推广条件。
% src:求解/问题3/结果/灵敏度.json:灵敏度_参数扰动.7.情景;灵敏度_参数扰动.8.情景;灵敏度_参数扰动.5.情景;灵敏度_参数扰动.6.情景;灵敏度_参数扰动.1.情景;灵敏度_参数扰动.2.情景;灵敏度_参数扰动.13.情景;灵敏度_参数扰动.14.情景;灵敏度_参数扰动.11.情景;灵敏度_参数扰动.12.情景
% src:交接/实验记录.json:564.现象
